\HMTNarrativeHeading{From recovering a state to transporting a collection}

Compatible recovery makes it possible to extend the exposition from
one state to collections of states. An enriched reading uniquely
determines its antecedents when its memory distinguishes the
elements that share the same visible value. The inverse then
restores the state on the image of that reading. Compatibility
with refinement assembles the finite inverses into a reconstruction
of complete histories.

Transporting collections requires preserving membership. A bijection
of histories takes each collection to its image and each element to its
corresponding element. Unions, intersections, and relative complements
are transported through these operations. The restriction of the
bijection to a collection preserves its cardinality. Memory theory thus
acquires a precise link with operations on sets of histories.

Measure also requires the transport of weights. A symmetry compatible
with the refinement tree preserves cylinder measure when it respects
emissions and their multiplicities. Naturality expresses the
agreement between transporting and truncating. The proofs that
follow make these relations explicit on the objects defined in
the article.

The specific study of the continuum hypothesis develops genealogical
closure and its hierarchy of sets. Its cardinal result is formulated
on that closure and its internal power set. The relation incorporated
by this part is the recovery of states and the transport of collections
that maintain the structure of the generated continuum. The cardinal
statement preserves its specialized development and its expressly
identified domain.

This articulation places the continuum within the conservation
argument. The complete history, the collection to which it belongs,
and the measure of its cylinder require related operations with
different contents. The article brings together their compatibility
conditions to show how the double projection is maintained in the
passage from an individual reading to the organization of collections.
The unity of the exposition resides in these maps and in the proofs
that preserve their relations.
